% ECONOMICS_PROBLEM: {"id": "I38-465", "title": "Targeting need versus potential impact", "jel_code": "I38", "source_id": "OP-0374"}
% !TeX program = xelatex
\documentclass[11pt,a4paper]{article}
\usepackage[margin=23mm]{geometry}
\usepackage{fontspec}
\setmainfont{Latin Modern Roman}
\usepackage{amsmath,amssymb,mathtools}
\usepackage{xcolor,enumitem,booktabs,longtable,array,xurl,hyperref}
\definecolor{muted}{HTML}{53636C}
\hypersetup{colorlinks=true,urlcolor=blue,linkcolor=blue}
\setlength{\parindent}{0pt}
\setlength{\parskip}{5pt}
\newcommand{\R}{\mathbb R}
\newcommand{\E}{\mathbb E}
\newcommand{\N}{\mathbb N}
\newcommand{\ind}{\mathbf 1}
\newcommand{\Var}{\operatorname{Var}}
\newcommand{\Cov}{\operatorname{Cov}}
\newcommand{\supp}{\operatorname{supp}}
\DeclareMathOperator*{\argmin}{arg\,min}
\DeclareMathOperator*{\argmax}{arg\,max}
\newcommand{\entrymeta}[3]{{\sffamily\small\color{muted}\textbf{\texttt{#1}}\quad|\quad #2\par #3\par}}
\newcommand{\Problemref}[1]{\texttt{#1}}
\newcommand{\JELref}[2]{\texttt{#2} (JEL \texttt{#1})}
\begin{document}
\section*{Targeting need versus potential impact}
\textbf{Problem ID:} \texttt{I38-465}\quad \textbf{JEL:} \texttt{I38}\par
% BEGIN ECONOMICS STATEMENT
\label{problem:OP-0374}
{\sffamily\small\color{muted}Original subject group: Development and poverty\par}
\label{definitions:problem:30007656}
\textbf{Audit classification:} Published research agenda. Published full text, author list and DOI checked. Other programs, contexts and longer horizons are explicit gaps; strategic-reporting and equity extensions are editorial.
\subsection*{Definitions and precise research target}
\textbf{Complete editorial problem.} Fix a target population, an assistance program, and a fixed per-household fiscal budget \(B\). Let \(X_i\) be baseline information legally and practically available for household \(i\), \(c_i\) the full delivery and targeting cost, and \(Y_i(d)\) a declared welfare outcome under eligibility \(d\in\{0,1\}\). Define need \(n_i\) with a predeclared deprivation measure and expected impact \(\tau(x)=E[Y_i(1)-Y_i(0)\mid X_i=x]\). A targeting rule \(\pi(x)\in[0,1]\) assigns eligibility probabilities. Define weighted gain \(g(x)=E[\omega(n_i)\{Y_i(1)-Y_i(0)\}\mid X_i=x]\). Characterize rules maximizing \(E[g(X_i)\pi(X_i)]\) subject to \(E[c_i\pi(X_i)]\le B\) and explicit minimum-coverage or equity constraints. The weight \(\omega\) is a declared social choice, not an estimated fact. Determine how the frontier changes when households manipulate reports, communities select recipients, or implementation creates exclusion and political capture. Randomized program assignment and independent measurement are needed to estimate impact; a deprivation predictor is not automatically an impact predictor. Validate rules in new settings and at scale, including administrative costs. Existing targeting studies answer important specific comparisons. Their agenda calls for other contexts, interventions, and longer horizons. The objective and strategic-reporting extensions here adapt that agenda; they are not a verbatim published formal problem.

\textbf{Definitions and admissible scope.} Let eligibility probabilities \(\pi(X)\) refer to assignment rather than take-up. Real resource cost \(c(X)>0\) includes targeting and delivery, and budget \(B\) is expected cost per eligible person. Need \(n\) and social weight \(\omega(n)\ge0\) are fixed in advance. If the welfare criterion is \(E[u(Y(\pi))]\), the proper gain is \(E[u(Y(1))-u(Y(0))\mid X]\), not generally \(\omega(n)E[Y(1)-Y(0)\mid X]\). The displayed linear weighted-gain objective is a separate explicit normative choice. Without interference, known gains and costs reduce to a knapsack or threshold allocation; uncertainty, external validity, strategic reports, equity constraints, and scalable implementation are the substantive extensions. For stochastic delivery costs write \(\bar c(x)=E[c_i\mid X_i=x]\); the base budget is \(E[\bar c(X)\pi(X)]\le B\), with \(\bar c>0\) and finite expectation. Need is a baseline variable and \(g(x)\), not \(\omega(n(x))\tau(x)\) unless need is fixed conditional on \(X\), is the correct linear weighted gain. An explicit group floor is \(E[\pi(X)\mid X\in G_k]\ge a_k\), for declared groups \(G_k\) of positive mass and \(a_k\in[0,1]\). For the base no-interference problem, the target is the set of policies attaining the supremum under these constraints; if gain and cost laws are uncertain, report value bounds or solve an explicitly declared minimax problem. Interference and report manipulation require new complete regime laws, as in the variants.
\subsection*{Variants and assumption changes}
\begin{enumerate}
\item Welfare curvature and horizon (source-motivated): compare targeting rules under declared \(u\) and discounted multiyear outcomes. Estimate \(E[u(Y_t(1))-u(Y_t(0))\mid X]\) rather than treating a mean outcome gain as exact expected utility.
\item Reporting and equity (editorial): permit applicants to choose reports with known costs and impose minimum coverage by deprivation group. Characterize policies using audited rather than self-reported covariates and estimate the welfare loss from manipulation.
\end{enumerate}
\subsection*{References and source audit}
\textbf{Support and access.} Published full text, author list and DOI checked. Other programs, contexts and longer horizons are explicit gaps; strategic-reporting and equity extensions are editorial. \textbf{Locator:} Section VI, pp. 1971--1972, caveats and final paragraph.
\begin{enumerate}\raggedright
\item\hypertarget{definitions:source:30007656:1}{} Johannes Haushofer; Paul Niehaus; Carlos Paramo; Edward Miguel; Michael W. Walker. 2025. \emph{Targeting Impact versus Deprivation}. Section VI, Conclusion, printed pp. 1971–1972, first two caveats and final paragraph.
\url{https://emiguel.econ.berkeley.edu/wordpress/wp-content/uploads/2026/01/haushofer-et-al-2025-targeting-impact-versus-deprivation.pdf}.
DOI: \url{https://doi.org/10.1257/aer.20221650}.
\end{enumerate}

\subsection*{Common conventions: Source mathematical conventions}

These conventions apply to each entry unless explicitly replaced.

\textbf{Models and identification.} A primitive model \(m\) specifies all probability laws, feasible choices, information, equilibrium and measurement rules used by its target. The admissible class \(\mathcal M\) is fixed independently of outcomes. If \(P_O(m)\) is its observable probability law and \(\tau(m)\) the target, its identified set at observable law \(P\) is
\[
 \mathcal I_\tau(P)=\{\tau(m):m\in\mathcal M,\ P_O(m)=P\}.
\]
Point identification means that this set is a singleton. An empty set signals incompatibility of data and assumptions. Coordinate bounds need not describe the joint set. Bounds are sharp only if every retained target is attainable or arbitrarily approachable by compatible models. Finite bounds require explicit restrictions on unobserved quantities; unrestricted confounding, hidden wealth, extrapolation or arbitrary equilibrium selection can yield uninformative sets.

\textbf{Causal interventions.} A potential outcome \(Y_i(p)\) is the outcome under a complete policy regime \(p\), including timing, financing, information and market spillovers. The target population and units are fixed before treatment. With interference, use the complete assignment vector or a justified exposure mapping; individual no-interference notation alone cannot identify an economy-wide intervention. A difference of mean potential outcomes does not require the cross-world joint distribution. Conditioning on post-treatment migration, survival, enrollment or employment changes the estimand and needs separate justification.

\textbf{Statistical statements.} An estimator, confidence set, forecast or test specifies a sampling law, model class and loss/coverage criterion. Uniform \((1-\alpha)\)-coverage means \(\inf_{m\in\mathcal M}\Pr_m\{\tau(m)\in C_n\}\geq1-\alpha\), or an explicitly asymptotic version. The whole parameter space trivially has coverage, so informativeness requires a stated width, risk or power objective. A forecasting improvement is relative to a named benchmark, information set and evaluation population.

\textbf{Equilibrium and optimization.} An equilibrium comprises feasible plans, optimality, beliefs where needed, budget constraints and all market/resource clearing. The primitives or a stated benchmark define each correspondence \(\mathcal E(p;m)\). Nonemptiness is assumed or proved, never inferred just from compact policy sets. A selected-equilibrium target uses a declared selection rule; a robust target quantifies over all permitted equilibria. Use suprema or infima unless attainment is established. An \(\varepsilon\)-optimal policy is feasible and within \(\varepsilon\) of the finite optimum value. Report an empty feasible set separately.

\textbf{Welfare and accounting.} Normative utility, interpersonal normalization, social weights, numeraire, discounting and the use of public receipts are primitives. Behavioral utility need not coincide with the welfare criterion. Household welfare includes ownership income and rents once; adding profits again to compensating variation generally double-counts them. A monetary equivalent is defined by an explicit transfer and price convention, with monotonicity and range conditions. Average effects, mechanism contributions, transfers and welfare losses are different targets. Infinite sums require convergence and dynamic feasible plans require terminal or no-Ponzi conditions.

\textbf{Source versions.} A working-paper date, revision date and journal publication year are distinguished. A DOI for a journal version is not automatically the DOI of an earlier manuscript. Locators refer to the stated version and printed pages where specified. A metadata-only check does not verify a claim about a full-text conclusion. Every entry's source subsection narrows the provenance claim accordingly.
% END ECONOMICS STATEMENT
\end{document}
